\documentclass[11pt,a4paper]{article}
\usepackage[utf8]{inputenc}
\usepackage[vietnamese]{babel}
\usepackage{amsmath,amssymb}
\usepackage{geometry}
\geometry{a4paper, total={180mm,260mm}, left=15mm, right=15mm, top=18mm, bottom=18mm}

\usepackage{tikz}
\usetikzlibrary{shapes.geometric, arrows.meta, positioning, calc}
\tikzset{arr/.style={-{Stealth[scale=0.8]}}}
\usepackage{tcolorbox}
\tcbuselibrary{skins,breakable}
\usepackage{fontawesome}
\usepackage{fancyhdr}
\usepackage{enumitem}

\pagestyle{fancy}
\fancyhf{}
\renewcommand{\headrulewidth}{0.5pt}
\renewcommand{\footrulewidth}{0.5pt}
\fancyhead[L]{\small\bfseries\color{violet!85!black}\faBook\ SỔ TAY KIẾN THỨC TOÁN 2}
\fancyhead[R]{\small\bfseries\color{violet!85!black}Lớp Toán Thầy Lĩnh -- loptoanthaylinh.vn}
\fancyfoot[C]{\small\color{gray}-- Sổ tay kiến thức: Bài 08. Phép cộng có nhớ trong phạm vi 20 --}

\setlength{\parindent}{0pt}
\setlength{\parskip}{0pt}

\begin{document}

% --- BANNER BÀI HỌC ---
\begin{tcolorbox}[
  enhanced,
  colback=violet!8!white,
  colframe=violet!70!black,
  arc=3mm, boxrule=1.2pt,
  left=3mm, right=3mm, top=2mm, bottom=2mm,
  halign=center
]
  {\LARGE\bfseries\color{violet!90!black}\faBookmark\ BÀI 08: PHÉP CỘNG (CÓ NHỚ) TRONG PHẠM VI 20}\\[1mm]
  {\itshape\footnotesize\color{black!75} Phương pháp tách số làm tròn 10 để tính nhẩm siêu tốc các dạng 9+, 8+, 7+, 6+ trên 1 trang A4}
\end{tcolorbox}

\vspace{-2mm}

\noindent
\begin{minipage}[t]{0.49\textwidth}\vspace{0pt}
\begin{tcolorbox}[
  enhanced,
  colback=violet!4!white,
  colframe=violet!75!black,
  arc=2mm, boxrule=0.8pt,
  left=2mm, right=2mm, top=1.8mm, bottom=1.8mm,
  title={\footnotesize\bfseries\color{white}\faKey\ 1. Phương pháp làm cho tròn 10}
]
{\footnotesize
\begin{itemize}\setlength{\itemsep}{1pt}\setlength{\leftskip}{-2.5mm}
  \item \textbf{Quy tắc chung:} Tách số thứ hai thành hai phần: một phần ghép với số thứ nhất cho đủ \textbf{10}, rồi cộng phần còn lại với 10.
  \item \textbf{Bảng tách theo số thứ nhất:}
  \begin{itemize}\setlength{\itemsep}{0.5pt}\setlength{\leftskip}{-2.5mm}
    \item \textbf{Dạng 9 + ...:} Tách $1$ ở số sau để được $9 + 1 = 10$.
    \item \textbf{Dạng 8 + ...:} Tách $2$ ở số sau để được $8 + 2 = 10$.
    \item \textbf{Dạng 7 + ...:} Tách $3$ ở số sau để được $7 + 3 = 10$.
    \item \textbf{Dạng 6 + ...:} Tách $4$ ở số sau để được $6 + 4 = 10$.
  \end{itemize}
\end{itemize}
}
\end{tcolorbox}
\end{minipage}
\hfill
\begin{minipage}[t]{0.49\textwidth}\vspace{0pt}
\begin{tcolorbox}[
  enhanced,
  colback=teal!4!white,
  colframe=teal!75!black,
  arc=2mm, boxrule=0.8pt,
  left=2mm, right=2mm, top=1.8mm, bottom=1.8mm,
  title={\footnotesize\bfseries\color{white}\faSliders\ 2. Ví dụ minh họa tách số}
]
{\footnotesize
\begin{itemize}\setlength{\itemsep}{1pt}\setlength{\leftskip}{-2.5mm}
  \item \textbf{Phép tính $9 + 5$:}
  Tách $5 = 1 + 4$.\\
  $9 + 5 = 9 + 1 + 4 = 10 + 4 = \mathbf{14}$.
  \item \textbf{Phép tính $8 + 5$:}
  Tách $5 = 2 + 3$.\\
  $8 + 5 = 8 + 2 + 3 = 10 + 3 = \mathbf{13}$.
  \item \textbf{Phép tính $7 + 6$:}
  Tách $6 = 3 + 3$.\\
  $7 + 6 = 7 + 3 + 3 = 10 + 3 = \mathbf{13}$.
  \item \textbf{Phép tính $6 + 7$:}
  Tách $7 = 4 + 3 \implies 6 + 4 + 3 = \mathbf{13}$.
\end{itemize}
}
\end{tcolorbox}
\end{minipage}

\vspace{-1.5mm}
\begin{tcolorbox}[
  enhanced,
  colback=orange!4!white,
  colframe=orange!80!black,
  arc=2mm, boxrule=0.8pt,
  left=2mm, right=2mm, top=1.5mm, bottom=1.5mm,
  title={\footnotesize\bfseries\color{white}\faSitemap\ 3. Sơ đồ tư duy tách số gộp 10}
]
\centering
\begin{tikzpicture}[x=1cm, y=0.9cm]
  \node[rectangle, rounded corners=2pt, draw=violet!80!black, fill=violet!15, font=\footnotesize\bfseries] (n1) at (0,0) {$9$};
  \node[font=\footnotesize\bfseries] at (1,0) {$+$};
  \node[rectangle, rounded corners=2pt, draw=teal!80!black, fill=teal!15, font=\footnotesize\bfseries] (n2) at (2,0) {$5$};

  \node[circle, draw=orange!80!black, fill=orange!20, inner sep=2pt, font=\scriptsize\bfseries] (p1) at (1.5,-1) {$1$};
  \node[circle, draw=orange!80!black, fill=orange!20, inner sep=2pt, font=\scriptsize\bfseries] (p2) at (2.5,-1) {$4$};

  \draw[->, thick, teal!70!black] (n2) -- (p1);
  \draw[->, thick, teal!70!black] (n2) -- (p2);

  \draw[arr, thick, violet!80!black, ->] (n1) to[bend right=30] node[below, font=\tiny\bfseries] {$9 + 1 = 10$} (p1);
  \node[rectangle, rounded corners=2pt, draw=green!80!black, fill=green!15, font=\scriptsize\bfseries] at (5,-0.5) {
    $10 + 4 = \mathbf{14}$
  };
\end{tikzpicture}
\end{tcolorbox}

\vspace{-1.5mm}
\noindent
\begin{minipage}[t]{0.49\textwidth}\vspace{0pt}
\begin{tcolorbox}[
  enhanced,
  colback=green!4!white,
  colframe=green!60!black,
  arc=2mm, boxrule=0.8pt,
  left=2mm, right=2mm, top=1.8mm, bottom=1.8mm,
  title={\footnotesize\bfseries\color{white}\faCheckCircle\ 4. Ví dụ thực hành chuẩn mực}
]
{\footnotesize
\textbf{Tính nhẩm nhanh:}\\
\begin{itemize}\setlength{\itemsep}{0.5pt}\setlength{\leftskip}{-2.5mm}
  \item $9 + 4 = 9 + 1 + 3 = \mathbf{13}$.
  \item $8 + 6 = 8 + 2 + 4 = \mathbf{14}$.
  \item $7 + 5 = 7 + 3 + 2 = \mathbf{12}$.
  \item $6 + 5 = 6 + 4 + 1 = \mathbf{11}$.
\end{itemize}
\textbf{Đổi chỗ số hạng:}\\
$4 + 9 = 9 + 4 = 13$ (đổi chỗ để cộng với số lớn trước cho nhanh).
}
\end{tcolorbox}
\end{minipage}
\hfill
\begin{minipage}[t]{0.49\textwidth}\vspace{0pt}
\begin{tcolorbox}[
  enhanced,
  colback=red!4!white,
  colframe=red!75!black,
  arc=2mm, boxrule=0.8pt,
  left=2mm, right=2mm, top=1.8mm, bottom=1.8mm,
  title={\footnotesize\bfseries\color{white}\faExclamationTriangle\ 5. Mẹo tránh bẫy trắc nghiệm}
]
{\footnotesize
\begin{itemize}\setlength{\itemsep}{1pt}\setlength{\leftskip}{-2.5mm}
  \item \textbf{Mẹo cộng với 9:}\\
  Hàng đơn vị của kết quả luôn \textbf{kém 1 đơn vị} so với số cộng cùng.\\
  VD: $9 + \mathbf{7} = 1\mathbf{6}$ ($7 - 1 = 6$); $9 + \mathbf{4} = 1\mathbf{3}$.
  \item \textbf{Mẹo cộng với 8:}\\
  Hàng đơn vị kết quả luôn \textbf{kém 2 đơn vị}: $8 + \mathbf{5} = 1\mathbf{3}$ ($5 - 2 = 3$).
  \item Luôn tách số bé để ghép đủ 10 với số lớn.
\end{itemize}
}
\end{tcolorbox}
\end{minipage}


\end{document}
